\section{Batch work inside a row-at-a-time executor}
\label{sec:customscan-batch-state-machine}

PostgreSQL's executor asks a child node for its next row. This follows the
iterator interface described in Volcano: the parent requests results one at
a time.\footnote{G.~Graefe, \emph{Volcano---An Extensible and Parallel Query
Evaluation System}, IEEE Transactions on Knowledge and Data Engineering
6(1), 1994, pp.~120--135, \url{https://doi.org/10.1109/69.273032}.}
A custom scan must
honor that interface even when the aligner's efficient work unit is a batch
of sequence pairs. BioSeqDB bridges the two sizes by retaining one batch and
returning its hits one at a time.

\begin{enumerate}
  \item \textbf{Collect input.} Read rows and apply supported ordinary
    conditions that do not depend on a hit. Retain the rows and sequence
    references needed for alignment.
  \item \textbf{Process a batch.} When the batch reaches its capacity, or the
    input ends, submit its sequence pairs to the selected alignment pipeline.
  \item \textbf{Return hits.} Select the next retained hit, supply its fields,
    and return one SQL row. Repeated requests advance through the same batch.
  \item \textbf{Reuse the space.} After all hits have been returned, release
    the completed batch and collect another one if input remains.
\end{enumerate}

\begin{figure}[htbp]
\centering
\begin{tikzpicture}[font=\footnotesize]
  \node[anchor=west,font=\bfseries] at (0,3.6) {Input batch};
  \foreach \y/\r in {2.65/A,1.75/B,.85/C}
    \node[scan input,minimum width=1.4cm] (row\r) at (.85,\y) {row \r};
  \node[scan hit,minimum width=2.5cm,minimum height=2.5cm] (align)
    at (4.0,1.75) {one batch\\alignment call};
  \foreach \r in {A,B,C}
    \draw[scan arrow] (row\r.east) -- (align.west);
  \node[anchor=west,font=\bfseries] at (6.25,3.6) {Retained hits};
  \node[scan hit,minimum width=.75cm] (a) at (6.75,2.65) {$A_0$};
  \node[scan hit,minimum width=.75cm] (b) at (7.75,2.05) {$A_1$};
  \node[text=black!60] at (7.25,1.40) {row B: no hit};
  \node[scan hit,minimum width=.75cm] (c) at (6.75,.85) {$C_0$};
  \draw[scan arrow] (align.east) -- (a.west);
  \draw[scan arrow] (align.east) -- (c.west);
  \node[anchor=west,font=\bfseries] at (9.4,3.6) {SQL rows};
  \node[scan output,minimum width=1.4cm] (out) at (10.25,2.65) {$A_0$};
  \draw[scan arrow,scanGreen] (a.east) -- (out.west);
  \node[align=center] at (10.25,1.5) {one per\\parent request};
  \node[anchor=west,text=black!65] at (0,-.1)
    {The next requests return $A_1$ and $C_0$ before another batch is collected.};
\end{tikzpicture}
\caption{Three input rows produce three hits with a different distribution:
two for A, none for B, and one for C. The parent still receives one row per
request. The aligner processes batches; SQL receives individual rows.}
\label{fig:scan-batch-pull}
\end{figure}

Scalar and vectorized table access converge before alignment submission. A
vectorized condition may first produce a mask of rows to keep, but surviving
rows still enter the same batch mechanism. PostgreSQL parallel execution is a
separate choice: the alignment owners are currently marked
\texttt{parallel\_safe=false}. The alignment pipeline may still use its own
workers or a GPU under its configuration.

\subsection{Row stability during a search}

PostgreSQL holds its current input row in temporary storage called a tuple
slot. The scan may still need that row for alignment after the slot is reused.
Another transaction may also update that row. The scan therefore separates
its saved values from the database view that it is allowed to read:

\begin{itemize}
  \item \textbf{Keep the input values.} A batch keeps ordinary fields in
    compact row copies in memory it owns, and keeps sequence references with
    their owners. A hit points to
    its batch entry, so several hits can return the same saved fields.
    Keeping a lazy reference does not require copying all sequence symbols.
  \item \textbf{Keep the same database view.} PostgreSQL uses a snapshot to
    decide which row versions a command can see. The sequence reader retains
    the scan's snapshot and checks a stored reference against the table,
    column, and source version visible through it. Stored source fragments
    are immutable. An update publishes a new version, while an
    older snapshot can still read its earlier version. Delaying a read does
    not switch it to the latest committed sequence.\footnote{PostgreSQL 17,
    \href{https://www.postgresql.org/docs/17/transaction-iso.html}
    {Transaction Isolation}.}
  \item \textbf{Separate internal and SQL identity.} The scan's internal
    row number locates source data; the SQL primary key identifies a logical
    row for the user. The batch-entry position selects saved values for a
    hit. None replaces the source-version and snapshot checks.
  \item \textbf{Keep the relations available.} Normal reads hold an
    \texttt{ACCESS SHARE} relation lock, and deferred source readers use the
    same mode. This prevents conflicting operations such as dropping or
    rewriting a relation while it is in use. Ordinary concurrent updates
    and deletes remain possible; the read does not lock each selected
    row.\footnote{PostgreSQL 17,
    \href{https://www.postgresql.org/docs/17/explicit-locking.html}
    {Explicit Locking}.}
\end{itemize}

A rescan discards the previous batch and hit fields while preserving the
scan's database view. Cleanup releases saved values before their source
readers and snapshots. Reuse across later queries also needs the result-cache
checks described below.

\subsection{Filter placement and early stopping}

Two kinds of condition require different inputs:

\begin{itemize}
  \item \textbf{A row condition} such as \texttt{t.id > 20} depends only on
    stored columns. Evaluating it before alignment avoids searching discarded
    rows.
  \item \textbf{A hit condition} needs an alignment result. The current
    implementation recognizes a narrow supported form such as
    \texttt{bioseq.edit(t.seq) < 10} and checks the edit count of each hit.
    This count is the sum of mismatched, inserted, and deleted symbols in
    that hit's alignment; each symbol counts once.
    This does not imply general support for arbitrary helper expressions in
    \texttt{WHERE}.
\end{itemize}

\noindent\begin{minipage}{\linewidth}
\begin{lstlisting}[style=bioseq-search-sql,caption={Filter input rows and then filter alignment hits.}]
SELECT t.id, bioseq.dp_score(t.seq) AS score
FROM targets AS t
WHERE t.id > 20
  AND mmsearch(t.seq, $1::bioseq.seq('dna5', 64))
  AND bioseq.edit(t.seq) < 10
LIMIT 5;
\end{lstlisting}
\end{minipage}

The identifier condition can avoid alignment for a row. The edit condition
selects hits after the required alignment information exists.

\texttt{LIMIT 5} lets the parent stop asking after five rows. It does not
guarantee that only five hits were computed: the current batch may already
be complete. Likewise, \texttt{ORDER BY score DESC LIMIT 5} does not by itself
give BioSeqDB a general top-$k$ alignment algorithm. Producing a correct
global score order can require reading all relevant hits.

\subsection{Initialization, rescans, and cleanup}

PostgreSQL defines four callbacks for the owner's execution
lifetime:\footnote{PostgreSQL 17,
\href{https://www.postgresql.org/docs/17/custom-scan-execution.html}
{Executing Custom Scans}.}

\begin{itemize}
  \item \texttt{BeginCustomScan} initializes execution state once.
  \item \texttt{ExecCustomScan} returns the next row, or reports the end.
  \item \texttt{ReScanCustomScan} restarts with the current input parameters.
  \item \texttt{EndCustomScan} releases execution resources.
\end{itemize}

Figure~\ref{fig:scan-lifecycle} shows the normal loop and the paths that need
special care.

\begin{figure}[htbp]
\centering
\begin{tikzpicture}[font=\footnotesize]
  \node[scan input,minimum width=2.6cm] (collect) at (2.1,3.1) {collect input};
  \node[scan hit,minimum width=2.6cm] (process) at (8.8,3.1) {process batch};
  \node[scan output,minimum width=2.6cm] (return) at (5.45,.5) {select next hit};
  \draw[scan arrow] (collect) -- node[above] {full batch or input ends} (process);
  \draw[scan arrow] (process.south) -- (return.east);
  \draw[scan arrow] (return.west) -- (collect.south);
  \node[align=center] at (6.5,2.0) {batch\\complete};
  \node[align=center] at (4.1,2.0) {no hit;\\more input};
  \draw[scan arrow] (return.south east) to[out=-35,in=-145,looseness=4]
    node[below,align=center] {return one hit;\\next row request} (return.south west);
  \node[scan box,minimum width=2.2cm] (init) at (2.1,5.0) {initialize};
  \draw[scan arrow] (init) -- (collect);
  \node[scan box,minimum width=2.2cm] (cleanup) at (11.25,.5) {release state};
  \draw[scan arrow] (return.east) -- node[pos=.55,above] {end or early stop} (cleanup.west);
  \draw[scan arrow,dashed,scanOrange] (process.east) to[out=0,in=90]
    node[right] {error} (cleanup.north);
  \draw[scan arrow,dashed,scanBlue,rounded corners=5pt] (return.west)
    -- (0,.5) -- (0,3.1) -- (collect.west);
  \node[anchor=west,align=left,text=scanBlue] at (.15,1.75)
    {rescan:\\clear hits,\\rebind inputs};
\end{tikzpicture}
\caption{One execution owns its batches and current hit. The ordinary loop
returns all hits of a batch before collecting another. A batch with no hits
continues to the next input batch or ends the scan. A normal end means that
no input rows or retained hits remain. A rescan starts with fresh inputs.
Cleanup is also required after early termination or an error.}
\label{fig:scan-lifecycle}
\end{figure}

\begin{itemize}
  \item \textbf{Initialize only execution state.} A prepared plan stores
    descriptions of inputs, not live alignment objects. Plain
    \texttt{EXPLAIN} follows a description-only initialization path and
    submits no alignment work.
  \item \textbf{Rescan with the new inputs.} A nested-loop join may supply a
    different query sequence on every outer row. The owner clears retained
    rows and hit cursors, rebinds the execution parameters, and restarts.
  \item \textbf{Release every owned resource.} Rows, sequence references,
    alignment buffers, and current-hit values have bounded lifetimes.
    PostgreSQL errors can unwind execution without ordinary C++ stack
    unwinding, so cleanup must also be registered with PostgreSQL's lifetime
    mechanisms. It cannot rely only on C++ destructors.
\end{itemize}
